\section{测试方案与结果}

\subsection{测试口径与证据范围}

滚球误差统一定义为$e_x=x_{\mathrm{ref}}-x$，并以全过程最大绝对误差
\begin{equation}
  E_{\max}=\max_{t\in\mathcal W}|e_x(t)|
  \label{eq:max-error}
\end{equation}
判定是否满足全程精度要求。相机最终标定为机械中心314像素，安全区间$[22,622]$像素，
两端分别约为$-11.6\,\si{cm}$和$+11.8\,\si{cm}$；正式换算采用\SI{25}{px/cm}的
保守比例。本文只统计保留下来的原始串口日志，不以程序设定值代替实测成绩。

\begin{table}[H]
  \centering
  \caption{报告所用证据及其边界}
  \label{tab:evidence}
  \small
  \begin{tabularx}{\linewidth}{p{3.2cm}Xp{3.1cm}}
    \toprule
    证据 & 可证明内容 & 不能证明内容 \\
    \midrule
    2026-08-01整车CSV（2304行） & 目标位置、整圈用时、分段球位误差、控制状态与丢球保护 & 多次重复性、正式裁判成绩、停车几何偏差 \\
    三位置视觉控制台记录 & 改进前中心、左端、右端的检测稳定性与帧率 & 最终滤波算法的漏检率 \\
    源代码与Keil构建记录 & 控制流程、保护逻辑、参数和可编译性 & 装车后的动态性能 \\
    MJPEG PC端单元测试 & 分片解析、客户端管理等软件逻辑 & 无线距离、现场清晰度、录像与回放 \\
    \bottomrule
  \end{tabularx}
\end{table}

\subsection{视觉标定与误检改进}

相机抬高、视野覆盖完整轨道后，先在旧中心标定（约320像素）下分别采集中心、左端和右端日志，
统计结果见表\ref{tab:vision-static}。中心位置因钢球反光和背景干扰出现5次无效帧，说明仅靠单帧圆检测
不足以保证稳定输出。随后加入运动预测选通、连续多帧获取确认和短时漏检保持，并将最终中心重新标定为314像素。

\begin{table}[H]
  \centering
  \caption{三位置视觉记录统计（算法改进前）}
  \label{tab:vision-static}
  \small
  \begin{tabular}{cccccc}
    \toprule
    位置 & 有效/无效帧 & 有效率 & 原始均值/px & 半径范围/px & 帧率/Hz \\
    \midrule
    中心 & 14/5 & 73.7\% & 319.86 & 9--18 & 95.8 \\
    左端 & 9/1 & 90.0\% & 42.89 & 11--17 & 98.0--98.2 \\
    右端 & 17/0 & 100.0\% & 617.29 & 16--19 & 98.4--98.7 \\
    \bottomrule
  \end{tabular}
\end{table}

纯逻辑测试覆盖了新目标三帧确认、预测保持、异常跳变拒绝、丢失后重新捕获和速度估计，全部通过。
该测试证明状态机实现符合设计，但不替代最终光照条件下的重复漏检率测量。

\subsection{整车闭环工程试验}

保留日志中的目标像素可由$x_{\mathrm{target}}=x_{\mathrm{ball}}+e_x$反推为506像素，
相对机械中心约为$+7.4\,\si{cm}$。从车辆进入运动状态到停车状态共\SI{28.00}{s}，
筛选滚球闭环有效、运动链路有效且车辆处于起步、直行或转弯状态的样本，共得到1400个控制点。
误差时序如图\ref{fig:r6-run}所示，统计结果见表\ref{tab:r6-result}和表\ref{tab:r6-segment}。

\begin{figure}[H]
  \centering
  \includegraphics[width=0.96\linewidth]{r6-engineering-run.png}
  \caption{偏置目标整车工程试验的球位误差与车辆运动状态}
  \label{fig:r6-run}
\end{figure}

\begin{table}[H]
  \centering
  \caption{整车工程试验总体结果}
  \label{tab:r6-result}
  \small
  \begin{tabular}{ccccccc}
    \toprule
    目标 & 圈时/s & MAE/cm & P95/cm & 最大误差/cm & $|e|\le1$ cm & 丢球事件 \\
    \midrule
    506 px & 28.00 & 0.78 & 1.52 & 1.76 & 69.3\% & 0 \\
    \bottomrule
  \end{tabular}
\end{table}

\begin{table}[H]
  \centering
  \caption{整车工程试验分运动阶段统计}
  \label{tab:r6-segment}
  \small
  \begin{tabular}{lccccc}
    \toprule
    阶段 & 样本数 & 等效时长/s & MAE/cm & 最大误差/cm & $|e|\le1$ cm \\
    \midrule
    起步/停车过渡 & 55 & 1.10 & 0.37 & 0.88 & 100.0\% \\
    直行 & 683 & 13.66 & 0.71 & 1.36 & 76.0\% \\
    右转 & 662 & 13.24 & 0.89 & 1.76 & 59.8\% \\
    \bottomrule
  \end{tabular}
\end{table}

该次试验的用时低于\SI{30}{s}，且没有发生从正常闭环到丢球保护的切换，证明视觉、两级串口和舵机控制在整圈运行中保持连续；然而最大误差\SI{1.76}{cm}超过\SI{1}{cm}，
因此只能判定“偏置目标整车闭环跑通且时间达到R6门槛”，不能判定R6整体通过。右转阶段的平均误差和超限比例均明显高于直行，说明主要矛盾已由静态摩擦转移到弯道惯性扰动、车体姿态变化及前馈相位匹配。

\subsection{各项任务完成状态}

\begin{table}[H]
  \centering
  \caption{赛题项目与现有证据对应关系}
  \label{tab:requirements}
  \small
  \begin{tabularx}{\linewidth}{cX>{\centering\arraybackslash}p{3.2cm}}
    \toprule
    项目 & 已完成工作 & 证据结论 \\
    \midrule
    R1 & K230双通道采集、MJPEG发送及PC端接收逻辑已实现 & 未留存现场完整录像与回放，不能判定通过 \\
    R2 & 循迹、起终点确认、TFT计时和主动停车已实现 & 未留存正式圈时与停车偏差重复记录 \\
    R3 & $314\rightarrow445\rightarrow190$像素自动序列、目标斜坡和稳定保持已实现 & 未留存完整时序曲线，不能报告正式成绩 \\
    R4 & A至B运行及滚球闭环功能已联调 & 未留存对应区间的独立原始记录 \\
    R5 & 中心目标整圈闭环功能已联调 & 未留存中心目标的完整重复记录 \\
    R6 & 506像素偏置目标完成一次整圈工程试验 & 28.00 s；最大误差1.76 cm，时间达标、精度未达标 \\
    \bottomrule
  \end{tabularx}
\end{table}

STM32工程最后一次Keil构建记录为0错误、0警告；MJPEG PC端7项单元测试全部通过。
这些结果用于说明软件可构建和数据链逻辑完整，不替代赛题要求的现场功能评价。

\subsection{误差来源与改进方向}

直行阶段仍有\SI{1.36}{cm}峰值，表明视觉比例误差、舵机死区、连杆间隙和滚球接触非线性仍然存在；
右转阶段峰值进一步增至\SI{1.76}{cm}，说明加速度及偏航前馈的方向、权重和相位还需在相同工况下做开关对照。
后续若继续优化，应固定机械结构与目标位置，分别记录“仅PID”“加速度前馈”“加速度加偏航前馈”三组重复试验，
以全过程最大误差而非平均误差选参，并保留每次原始CSV和同步视频。积分增益用于消除持续偏差，但不能替代弯道前馈；
过度增大积分会在入弯和出弯时积累并导致反向超调，因此应保持积分分离和抗饱和约束。
